\documentclass[paper=a4paper,fontsize=10.4pt]{jlreq}

%%% 余白設定
%\usepackage[top=3cm, bottom=3cm, left=2cm, right=2cm]{geometry}
\usepackage[driver=dvipdfm,truedimen]{geometry}
\geometry{left=15truemm,right=15truemm,top=12truemm,bottom=12truemm}

%%% フォント設定
\usepackage{luatexja-fontspec}
\setmainjfont{UDDigiKyokashoNP}[YokoFeatures={JFM=jlreq},TateFeatures={JFM=jlreqv}]
\setmainfont{UDDigiKyokashoNP}

\usepackage{amssymb,amsmath}
%\usepackage[dvipdfmx,hidelinks]{hyperref}
\usepackage[pdfencoding=auto]{hyperref}
\usepackage{tikz}

%%% 枠囲み
\usepackage{tcolorbox}
\tcbuselibrary{skins}

\usepackage{enumitem}
\setlist{leftmargin=5mm}
\renewcommand{\labelitemi}{$\triangleright$~}
\renewcommand{\labelitemii}{$\blacktriangleright$~}

\usepackage{multicol}

%%% section のフォント設定
%\usepackage{titlesec}
%\titleformat{\section}{\Large\bfseries}
%  {\thesection}{1em}
%  {}
%  [\titlerule]
%\titleformat*{\subsection}{\Large\bfseries}
%\titleformat*{\subsubsection}{\large\bfseries}
\ModifyHeading{section}{%
  font={\Large\bfseries},
  format={%
    \begin{tikzpicture}
      \useasboundingbox (0,0) rectangle (0,0);
      \fill[black!50] (0,-4pt) rectangle (8pt,12pt);
      \draw[black!50,line width=1pt] (0,-4pt) -- (\linewidth,-4pt);
    \end{tikzpicture}%
    \hspace{1\zw}#1#2
  },
  before_lines=0,
  %lines=2,
  after_lines=-.5,
}
\ModifyHeading{subsection}{font={\Large\bfseries},before_lines=0,after_lines=0}
\ModifyHeading{subsubsection}{font={\small\bfseries},before_lines=0,after_lines=.5}

\renewcommand{\emph}[1]{\textbf{#1}}

\pagestyle{empty}

\title{計算情報学研究室}
%\author{深川 大路 (Fukagawa, Daiji)}
\date{}



\begin{document}

\begin{minipage}[t]{.45\textwidth}
\noindent
{\LARGE\bfseries\rmfamily {\Huge ㉘} 計算情報学研究室}\\
\hspace{.4cm} Computational Informatics Laboratory
\end{minipage}
\begin{minipage}[t]{.45\textwidth}
\begin{tabular}{l@{ }l}
担当教員: & 深川 大路 (FUKAGAWA Daiji) \\
e-mail: & \texttt{dfukagaw@mail.doshisha.ac.jp} \\
研究室: & MK515 (教員)， MK211 (学生)
\end{tabular}
\end{minipage}

%\maketitle

\section{教員について}

専門分野は\emph{アルゴリズムと計算量の理論}，特に，\emph{離散データのマッチング}です。
ざっくりと説明すると，アルゴリズムとは問題の解き方であり，計算量はそのアルゴリズムの効率を表す尺度のひとつであり，離散データは，文字列や木などの単純な構造を持つデータから，グラフとよばれる複雑なネットワーク構造を持つデータを含むものを指します。

アルゴリズムを研究する最大の目的は，計算機 (computer) の処理を速くすることです。そのための手段として，アルゴリズムの計算量や問題の複雑さを数学的に解析したり，実際に計算機プログラムを実装したりといった作業を行います。

一方で，講義・研究・その他の用務上の必要からプログラムを書くことも多く，その関係から，\emph{情報システムの開発}にも大いに興味があります。最近はこちらに重点を置くことが多くなっています。利用するプログラミング言語は最近では Python と JavaScript が多く，他にも C/C++, Java, Ruby などを使います。

\section{卒業研究について}

卒業研究のテーマは，配属後なるべく自分自身で考えてもらいます。ただし何でもありではなく，教員とディスカッションを重ね，取り組むに値するかどうかを吟味して決定します。その過程において，しっかりと先行研究を調査してください。教員から批判的な意見も投げかけますので，やりたいこと・考えたこと・調べたことをきちんと説明してください。ここで手を抜くと後で苦労した挙句に行き詰まり，泣く泣くテーマを変えるということになりかねません。テーマの種を見極める作業は意外と重要です。
教員自身が抱えているテーマが自分の興味にマッチしていれば，それを選んでいただくこともオススメです。教員はモチベーションが上がります。ただし実際に卒業研究を行う皆さんにとってのモチベーションがないとお互いに困ります。

\begin{tcolorbox}[
colframe=gray, colback=white,
coltitle=white, colbacktitle=gray,
boxrule=0.8pt, arc=0mm,
outer arc=0mm,
sharp corners,
fonttitle=\bfseries,
enhanced,
attach boxed title to top left={xshift=5mm,yshift=-3mm},
boxed title style={frame hidden},
title=過去の卒業研究題目（過去4期分）
]
%\vspace{-.5cm}
\begin{multicols}{2}
\fontsize{8bp}{11pt}\selectfont
\subsubsection*{2024年度}
\begin{itemize}[leftmargin=3mm]
  \item 生物学的制約の実装による人間らしい挙動をするゲームAIの作成
  \item 機械学習によるサッカーの得点数予測
  \item 縮小版マンカラ・ベーシックの強解決
  \item カンディンスキー作品と拡散モデルの画像の色彩分析
  \item 最短経路探索における粘菌アルゴリズムの性能評価
  \item mini2048の完全解析
  \item 動的難易度調整を行う格闘ゲームAIへの生物学的制約の導入
  \item 骨格認識技術を用いた転倒時の特徴分析
  \item 人狼ゲームにおける人狼陣営の勝率向上を図るAIの提案
  \item 人狼ゲームにおける裏切り者を起点とするパワープレイ戦略の提案
\end{itemize}
\subsubsection*{2023年度}
\begin{itemize}[leftmargin=3mm]
  \item 遺伝的アルゴリズムを用いた文字ジグソーパズルの解法
  \item 選好順序の偏りを考慮した研究室配属のシミュレーション
  \item Sentence-BERTを用いた商品データのカテゴリー分類
  \item スポーツ動画における得点シーンの自動推定
  \item ぷよぷよの戦術についての統計的分析 －中盤戦に着目して－
  \item ゴールデンステート・ウォリアーズがNBA全体に与えた影響の考察
  \item 楽天市場サイトにおける商品の使用感の自動抽出
  \item Pongにおけるエージェントの生物的制約を加えた人間らしさの獲得
\end{itemize}
\subsubsection*{2022年度}
\begin{itemize}[leftmargin=3mm]
  \item チャンツを用いた英語リズム習得---英語習熟度とリズム感が学習効果に及ぼす影響---
  \item ハンドクラフト特徴量を考慮したファッションアイテムの類似度判定方法
  \item 複数の画像変換を組み合わせた複雑背景物体検出
  \item 解答自由度に着目した国語記述式問題の自動採点
  \item 音楽のBPMの違いが集中力へ与える影響---作業前の聴取に着目して---
  \item 阪神高速11号池田線における交通流の特徴分析および渋滞対策の検討
\end{itemize}
\subsubsection*{2021年度}
\begin{itemize}[leftmargin=3mm]
  \item 過去の着手をLSTMにより考慮するガイスターの駒色推定手法について
  \item 学生の履修履歴とシラバスに基づく履修選択支援システム
  \item Pommermanにおける協働設定が強化学習エージェントに与える影響について
\end{itemize}
\end{multicols}
\end{tcolorbox}

\clearpage
\section{卒業研究の進め方}

卒業研究は個人で行います。
方向性や進捗を確認できるよう，\emph{全体ミーティング（毎週）}と\emph{個別ミーティング（随時）}を活用してください。
人数が少ない年度は個別ミーティング中心，人数が多い年度は全体ミーティング中心となります。
たとえば2020年度は卒研生を3つのゼミ（アルゴリズム，機械学習，情報システム）に分割してゼミ単位でミーティングを行いました。
2021～2024年度は週2回の全体ミーティング中心です。
週4コマのゼミの中で勉強会（プログラミング，機械学習等）も実施しています。

研究室に配属が決まったら遠慮せずに研究室に来てください。
ゼミ・勉強会に参加することも可能です。
先輩の様子を見て話を聞いてみてください。卒業研究において必要なものは何か，考えてください。
見る・聞くだけでなく自分でやってみることも大切です。
\emph{積極性・主体性}を重視します。\emph{2年次生のうちから先輩の卒業研究を引き継ぎ，改善したものを学外で発表して受賞した先輩たち}もいます。ぜひチャレンジしてください。

\subsection*{3年次生}
\vspace{-.2cm}
\begin{itemize}[topsep=0pt]
  \item 必修科目ジョイントリサーチで，卒業研究に必要な知識・技術・姿勢を身に付けます。
  \item 日々更新される最新技術・専門知識も重要ですが，まずは基礎力を鍛えることを優先してください。
\end{itemize}

\subsection*{4年次生}
\begin{itemize}[topsep=0pt]
  \item 春学期：期末レポート（初稿6月初旬，最終稿7月初旬），コロキアム発表練習
  \item 秋学期：コロキアム発表，卒業論文（初稿10月下旬，最終稿11月下旬），試問会（1月下旬）
\end{itemize}

\section{教員コンタクト（配属希望者のための面接）}

\begin{itemize}[topsep=0pt]
  \item 詳しい日程等は，学部全体・コース全体の研究室配属関連資料を参照してください。
  \item まず電子メールでアポイントメントをとってください。
    少なくとも一通目は，面識のない相手に送るビジネス文書のつもりで書いてみてください（その能力があることを見せてください）。
    メールには日程調整を円滑に行うために必要な情報を含めてください。
  \item 面接は対面で実施します。やむを得ずオンラインを希望される場合は（余裕を持って）相談してください。
  \item 面接の際には，成績が確認できるものを持参してください（成績表全体を印刷したものが望ましいです）。
  \item 教員コンタクト時に確認する内容
    \begin{itemize}
      \item 卒業研究に対する準備状況・今後の計画について自身の言葉で説明できるようまとめておいてください。
      \item 〔準備状況〕これまでに履修した科目や個人的に学んだ事柄
      \item 〔今後の計画〕時系列に沿って（どの時期に何をするか）可能な限り具体的に
      \item 研究テーマとこれまで履修した科目との関連性についてランダムに尋ねます。簡潔に（１科目あたり３０秒～１分程度で）説明できるように準備しておいてください。
    \end{itemize}
\end{itemize}

\section{配属希望者へのメッセージ}

最も重要なことはマッチングです。面談での受け答えを重視します。とはいえ緊張し過ぎて自分を出せないのも困ります。自信を持ってあなたのことを十分に説明してください。

必要であれば何でもやってみる気持ちで広く勉強をしてください。そのうえで，ここが重要だというところは，とことん深く掘り下げてください。理解した「つもり」は危険です。
今はまだできないこと，分からないことであっても，日々の努力の積み重ねで乗り越えて，
できること，分かることを貪欲に広げましょう。

「研究対象への深い興味」「研究の遂行に必要な知識を探し求める姿勢」
「卒業研究に真剣に取り組む誠実さ」「日常生活を楽しめる前向きな姿勢」
を備えた人を募集します。真剣に研究をしたい人を力の限り応援します！

%
%
\end{document}
